\Askhsh{1754}{4}
\begin{ylh}{1}
    \item 3.4. 3ο Κριτήριο ισότητας τριγώνων
    \item 3.6. Κριτήρια ισότητας ορθογώνιων τριγώνων
    \item 5.8. Το ορθόκεντρο τριγώνου.
\end{ylh}
Δίνεται ισοσκελές τρίγωνο $\mathrm{A}\mathrm{B}\Gamma$ και το ύψος του $\mathrm{A}\Delta$. Στο $\mathrm{A}\Delta$ θεωρούμε σημείο $\Theta$ τέτοιο ώστε $\Theta\mathrm{A} = \Theta\mathrm{B}$. Έστω ότι $\mathrm{E}$ είναι το σημείο τομής της $\mathrm{B}\Theta$ με την $\mathrm{A}\Gamma$. Φέρνουμε την $\mathrm{A}\mathrm{H}$ κάθετη στην $\mathrm{B}\mathrm{E}$, η οποία τέμνει την πλευρά $\mathrm{B}\Gamma$ στο $\mathrm{Z}$.
\begin{alist}
    \item Να αποδείξετε ότι:
    \begin{rlist}
        \item Τα τρίγωνα $\Theta\Delta\mathrm{B}$ και $\Theta\mathrm{H}\mathrm{A}$ είναι ίσα. \monades{6}
        \item $\Delta\Theta = \Theta\mathrm{H}$. \monades{6}
        \item Η ευθεία $\Delta\Theta$ είναι μεσοκάθετος του τμήματος $\mathrm{A}\mathrm{B}$. \monades{6}
    \end{rlist}
    \item Ποιο από τα σημεία του σχήματος είναι το ορθόκεντρο του τριγώνου $\mathrm{A}\Theta\mathrm{B}$; Να δικαιολογήσετε την απάντησή σας. \monades{7}
\end{alist}
